\documentclass[12pt,a4paper]{ctexart}

% ===== 版式 =====
\usepackage[top=2.5cm,bottom=2.5cm,left=2.6cm,right=2.6cm]{geometry}
\usepackage{amsmath,amssymb}
\usepackage{booktabs,multirow,longtable,array}
\usepackage{graphicx}
\usepackage[labelsep=quad,font=small]{caption}
\usepackage{subcaption}
\usepackage{enumitem}
\usepackage[colorlinks,linkcolor=black,citecolor=black,urlcolor=black]{hyperref}

% 图路径：正式稿把图放 图集/Mx/，这里指向仓库内相对路径
\graphicspath{{../../图集/}}

\title{【题目名称】}
\author{队伍编号：\underline{\hspace{3cm}}}
\date{}

\begin{document}

% ===== 封面/标题 =====
\maketitle

% ===== 摘要（数字、专有名词、"针对问题一二三"加粗）=====
\begin{abstract}
\noindent 本文针对【问题背景一句话】，建立了【模型名称】。

针对问题一，\textbf{【方法】}……，得到\textbf{【关键数字】}……

针对问题二，\textbf{【方法】}……，得到\textbf{【关键数字】}……

针对问题三，\textbf{【方法】}……，得到\textbf{【关键数字】}……

本文模型经过误差分析、灵敏度检验与稳健性检验，结果【一句话结论】。

\vspace{0.6em}
\noindent\textbf{关键词：}关键词一；关键词二；关键词三；关键词四；关键词五；关键词六
\end{abstract}

\newpage

% ===== 正文：各问统一格式 问题描述 → 模型建立 → 算法求解 → 求解结果 =====
\section{问题重述}
\subsection{问题背景}
\subsection{问题提出}

\section{问题分析}
\subsection{问题一的分析}
\subsection{问题二的分析}

\section{模型假设}
\begin{enumerate}[label=(\arabic*)]
    \item 假设一：……
    \item 假设二：……
\end{enumerate}

\section{符号说明}
\begin{table}[htbp]
    \centering
    \caption{符号说明}
    \begin{tabular}{lll}
        \toprule
        符号 & 含义 & 单位 \\
        \midrule
        $x_{i}$ & 决策变量 & 个 \\
        $\beta$ & 回归系数 & -- \\
        \bottomrule
    \end{tabular}
\end{table}

\section{问题一模型的建立与求解}
\subsection{问题描述}
\subsection{模型建立}
\begin{equation}
    \min_{x} \; f(x) = \sum_{i=1}^{n} \left( y_i - \hat{y}_i \right)^2
    \quad \text{s.t.} \quad
    \begin{cases}
        A x \le b, \\
        x \ge 0 .
    \end{cases}
\end{equation}

\subsection{算法求解}
\subsection{求解结果}
% 插图：只插 图集/Mx/ 里的图，不重画
\begin{figure}[htbp]
    \centering
    % \includegraphics[width=0.75\linewidth]{M1/图M1-01-示例.png}
    \fbox{\parbox[c][4cm][c]{0.6\linewidth}{\centering 图片占位：引用 图集/Mx/图Mx-NN-短名.png}}
    \caption{【图题】}
\end{figure}

\section{模型检验}
\subsection{误差分析}
\subsection{灵敏度分析}
\subsection{稳健性分析}

\section{模型评价与推广}
\subsection{模型优点}
\subsection{模型缺点}
\subsection{模型推广}

% ===== 参考文献：6–7 篇，英文 ≤2 篇（专家十条第 1 条）=====
\begin{thebibliography}{9}
\bibitem{ref1} 姜启源, 谢金星, 叶俊. 数学模型（第五版）. 北京: 高等教育出版社, 2018.
\bibitem{ref2} 司守奎, 孙玺菁. 数学建模算法与应用（第三版）. 北京: 国防工业出版社, 2021.
\end{thebibliography}

\end{document}
